\documentclass[12pt]{article}

\usepackage[a4paper, margin=3cm]{geometry}
\usepackage{times}
\usepackage{parskip}
\usepackage{microtype}
\usepackage{hyperref}

\hypersetup{colorlinks=true, urlcolor=black, linkcolor=black}

\pagestyle{empty}

\begin{document}

\noindent
Madoka Masui\\
4-7-14 Zaimokuza, Kamakura, 248-0013 Japan\\
lucaremm@gmail.com

\bigskip
\noindent
\today

\bigskip
\noindent
Dear Editor,

\bigskip
I am pleased to submit the manuscript entitled \textit{``Irreversibility-Constrained
Governance: Preserving Human Responsibility under Finite Oversight Capacity''} for
consideration in \textit{Ethics and Information Technology}.

The manuscript addresses a structural problem emerging in AI-mediated institutional
environments: under conditions of continuous updating and adaptive optimization, human
responsibility may erode not through misattribution, but through the progressive depletion
of oversight capacity prior to irreversible outcomes.

While existing governance frameworks such as Meaningful Human Control clarify alignment
and traceability conditions, they do not model the dynamic evolution of institutional
oversight capacity itself. The proposed framework, Irreversibility-Constrained Governance
(ICG), introduces a capacity-based model of institutional irreversibility and formalizes
a velocity constraint regulating the rate at which systems approach configurations in which
evaluative and corrective judgment become structurally ineffective.

The contribution of the manuscript is threefold:

\begin{enumerate}
  \item It distinguishes biological irreversibility from institutional irreversibility as
        analytically separate but structurally related phenomena.
  \item It models oversight capacity as a dynamic variable subject to asymmetric depletion
        and recovery, incorporating institutional hysteresis.
  \item It proposes structured freeze mechanisms as recovery operations that preserve the
        evaluative space necessary for responsible judgment.
\end{enumerate}

The framework is instantiated in a high-risk ICU sepsis decision-support context to
demonstrate that the structural dynamics of capacity depletion and recovery are most
consequential where biological irreversibility represents an absolute downstream threshold.

The manuscript does not argue against technological optimization. Rather, it proposes
that governance must regulate the trajectory of system evolution under conditions of
finite human capacity. Its objective is to preserve the structural conditions under which
human responsibility remains operational rather than merely formal.

This manuscript has not been published elsewhere and is not under consideration by any
other journal. There are no conflicts of interest to declare. This work received no
external funding and was conducted without institutional affiliation.

Thank you for your consideration.

\bigskip
\noindent
Sincerely,

\vspace{1.2cm}
\noindent
Madoka Masui\\
Independent Researcher\\
4-7-14 Zaimokuza, Kamakura, 248-0013 Japan\\
lucaremm@gmail.com

\end{document}
